% ===========================================================================
%  §4  TOWN: Think Only When Needed
%  method_final.tex — NeurIPS 2026
%  Loaded via % ===========================================================================
%  §4  TOWN: Think Only When Needed
%  method_final.tex — NeurIPS 2026
%  Loaded via % ===========================================================================
%  §4  TOWN: Think Only When Needed
%  method_final.tex — NeurIPS 2026
%  Loaded via % ===========================================================================
%  §4  TOWN: Think Only When Needed
%  method_final.tex — NeurIPS 2026
%  Loaded via \input{sections/method_final}
% ===========================================================================

The empirical findings in \S\ref{sec:thinking-tax} suggest a clear strategy:
\textbf{default to non-thinking mode and reserve thinking for genuinely
hard problems}.
We formalize this intuition as \textsc{Town} (\textbf{T}hink
\textbf{O}nly \textbf{W}hen \textbf{N}eeded), a training-free,
two-stage inference cascade that requires no access to model internals
and no auxiliary models.

% ---------------------------------------------------------------------------
\subsection{Key Insight: Early Stopping as a Difficulty Classifier}
\label{sec:method:insight}
% ---------------------------------------------------------------------------

Our analysis reveals that non-thinking mode's early-stopping behavior
provides a natural binary classifier for problem difficulty:

\begin{observation}
\label{obs:early-stop}
For Qwen3-8B with \texttt{nothink@256} on GSM8K ($n{=}1{,}319$):
\begin{itemize}[nosep,leftmargin=*]
    \item \textbf{88.8\% of samples stop early} (avg 133 tokens),
      achieving \textbf{94.4\%} accuracy among early-stopped samples
      (87.5\% overall).
    \item \textbf{11.2\% of samples hit the budget} (use all 256 tokens),
      indicating the model struggled.
\end{itemize}
\end{observation}

This dichotomy provides a zero-cost routing signal: if the model finishes
quickly in non-thinking mode, the answer is likely correct; if it
exhausts the budget, the problem is likely hard and may benefit from
extended reasoning.
Unlike soft confidence scores or logit-based uncertainty measures, this
signal is \emph{binary}, \emph{endogenous}, and \emph{free}---it emerges
as a byproduct of the generation process itself
(Finding~2, \S\ref{sec:finding:natural-stop}).

% ---------------------------------------------------------------------------
\subsection{The \textsc{Town} Cascade}
\label{sec:method:cascade}
% ---------------------------------------------------------------------------

\textsc{Town} operates in two stages, summarized in
Algorithm~\ref{alg:town} and Figure~\ref{fig:town-pipeline}.

\begin{algorithm}[t]
\caption{\textsc{Town}: Think Only When Needed}
\label{alg:town}
\begin{algorithmic}[1]
\REQUIRE Question $q$, model $\mathcal{M}$,
  probe budget $B_1$ (non-thinking),
  fallback budget $B_2$ (thinking, $B_2 > B_1$)
\ENSURE Answer $\hat{a}$, total tokens consumed $T$

\medskip
\STATE \textbf{// Stage 1: Non-thinking probe}
\STATE $y_1 \leftarrow \mathcal{M}_{\textrm{nt}}(q,\; B_1)$
  \hfill\COMMENT{Generate without thinking}
\STATE $T \leftarrow |y_1|$

\medskip
\IF{$|y_1| < B_1$}
  \hfill\COMMENT{Natural stop $\Rightarrow$ high confidence}
  \STATE $\hat{a} \leftarrow a(y_1)$
  \hfill\COMMENT{Accept non-thinking answer}
\ELSE
  \hfill\COMMENT{Budget hit $\Rightarrow$ low confidence}
  \STATE \textbf{// Stage 2: Thinking fallback}
  \STATE $y_2 \leftarrow \mathcal{M}_{\textrm{th}}(q,\; B_2)$
    \hfill\COMMENT{Generate with chain-of-thought}
  \STATE $T \leftarrow T + |y_2|$
  \STATE $\hat{a} \leftarrow a(y_2)$
  \hfill\COMMENT{Use thinking answer}
\ENDIF

\medskip
\RETURN $\hat{a},\; T$
\end{algorithmic}
\end{algorithm}

\input{sections/fig_town_pipeline}

\paragraph{Stage 1: Non-thinking probe.}
We run the input through non-thinking mode with a moderate budget
$B_1$ (default: 256 tokens).
If the model produces a complete answer before exhausting $B_1$, we
accept it immediately.
This stage handles the vast majority of inputs---88.8\% on GSM8K---at
minimal cost (${\sim}$133 tokens on average for early-stopped samples), exploiting Finding~1's
observation that non-thinking mode is more token-efficient than thinking
mode at matched budgets.

\paragraph{Stage 2: Thinking fallback.}
For the ${\sim}$11.2\% of inputs that exhaust the Stage~1 budget---indicating
the model could not produce a confident answer quickly---we route to
thinking mode with an extended budget $B_2$ (default: 512 tokens).
These are precisely the ``hard'' problems where step-by-step reasoning
chains are most likely to provide genuine benefit.
The total cost for a fallback sample is $B_1 + |y_2|$, since the
probe consumed its full budget.

% ---------------------------------------------------------------------------
\subsection{Cost and Accuracy Analysis}
\label{sec:method:analysis}
% ---------------------------------------------------------------------------

\paragraph{Expected token cost.}
Let $p$ denote the early-stop rate in Stage~1,
$\bar{t}_1$ the average tokens for early-stopped samples, and
$\bar{t}_2$ the average thinking-mode cost for fallback samples.
The expected token cost per question is:
\begin{equation}
  \label{eq:town-cost}
  \mathbb{E}[T] \;=\; p \cdot \bar{t}_1
    \;+\; (1 - p) \cdot \big(B_1 + \bar{t}_2\big).
\end{equation}
With our GSM8K parameters ($p{=}0.888$, $\bar{t}_1{\approx}133$,
$B_1{=}256$, $\bar{t}_2{\approx}469$) on the full test set ($n{=}1{,}319$):
\begin{equation}
  \label{eq:town-cost-concrete}
  \mathbb{E}[T]
    \;\approx\; 0.888 \times 133 \;+\; 0.112 \times (256 + 469)
    \;\approx\; 118 + 81
    \;=\; \mathbf{\sim\!199\text{ tokens}}.
\end{equation}
This is \textbf{2.4$\times$} cheaper than \texttt{think@512} alone
(460~tokens), representing a \textbf{57\%} reduction in average token
consumption.  Note that the Stage~1 probe budget $B_1$ is ``spent''
even for routed samples---it serves as the cost of difficulty
classification.

\paragraph{Projected accuracy.}
\textsc{Town}'s accuracy decomposes as:
\begin{equation}
  \label{eq:town-acc}
  \textrm{Acc}_{\textsc{Town}}
    \;=\; p \cdot \textrm{Acc}_{\textrm{nt}}^{\textrm{early}}
    \;+\; (1 - p) \cdot \textrm{Acc}_{\textrm{th}}^{\textrm{hard}},
\end{equation}
where $\textrm{Acc}_{\textrm{nt}}^{\textrm{early}}$ is the
accuracy of early-stop non-thinking samples and
$\textrm{Acc}_{\textrm{th}}^{\textrm{hard}}$ is the thinking-mode
accuracy on the hard subset routed to Stage~2.
On the full GSM8K test set ($n{=}1{,}319$),
\textsc{Town} achieves \textbf{90.9\%} accuracy at only \textbf{199
average tokens} (including the Stage~1 probe cost for routed samples):
\begin{equation}
  \textrm{Acc}_{\textsc{Town}}
    \;=\; \underbrace{0.888 \times 94.4\%}_{\text{Stage 1: early-stop}}
    \;+\; \underbrace{0.112 \times 62.8\%}_{\text{Stage 2: thinking fallback}}
    \;=\; 83.8\% + 7.0\%
    \;=\; \mathbf{90.9\%}.
\end{equation}
This exceeds \texttt{think@512} (65.2\%) by \textbf{+25.7\,pp}
while using \textbf{2.4$\times$ fewer tokens}, and surpasses
\texttt{nothink@256} (87.5\%) by \textbf{+3.4\,pp} by recovering
hard cases through Stage~2.
Full evaluation results are in \S\ref{sec:town-eval}.

% ---------------------------------------------------------------------------
\subsection{Comparison with Alternative Strategies}
\label{sec:method:comparison}
% ---------------------------------------------------------------------------

\paragraph{vs.\ \texttt{think@512} only.}
\textsc{Town} achieves substantially higher accuracy at
${\sim}$2.4$\times$ fewer tokens.
The advantage comes from avoiding thinking mode's truncation waste on the
88.8\% of problems where it is unnecessary.

\paragraph{vs.\ \texttt{nothink@256} only.}
\textsc{Town} adds a thinking fallback for the $\sim$11.2\% of hard
problems, recovering accuracy on cases where direct answering
fails.
The additional cost is modest: $\sim$53 extra tokens on average
over the entire dataset (199 vs.\ 146 for \texttt{nothink@256}).

\paragraph{vs.\ Self-consistency~\cite{wang2023selfconsistency}.}
SC@$K$ generates $K$ full-budget reasoning traces and majority-votes,
consuming $K \times B$ tokens unconditionally.
\textsc{Town} uses a \emph{single} non-thinking probe and a binary
routing decision, yielding lower latency (one generation call for easy
questions vs.\ $K$), lower token cost ($\sim$146 tokens for easy
questions vs.\ $K \times B$), and a simpler implementation (no consensus
computation, no threshold tuning).

\paragraph{vs.\ Oracle routing.}
An oracle that knows each problem's difficulty \emph{a priori}
achieves 68.2\% accuracy at 401 average tokens
(Table~\ref{tab:efficiency-frontier}).
\textsc{Town}'s early-stop signal approximates this oracle using a
training-free heuristic, with the gap determined primarily by the
signal's precision on the routed subset.

\paragraph{Hyperparameters.}
\textsc{Town} introduces exactly \emph{two} hyperparameters beyond the
base model: $B_1$ (probe budget) and $B_2$ (thinking fallback budget).
Both have intuitive interpretations: $B_1$ should be large enough for
the model to produce a complete non-thinking answer on most easy
questions (256 tokens suffices for GSM8K), and $B_2$ should be large
enough for a complete reasoning chain on hard questions (512+ for
GSM8K).
No routing thresholds, consensus weights, or learned parameters are
required.

% ===========================================================================

The empirical findings in \S\ref{sec:thinking-tax} suggest a clear strategy:
\textbf{default to non-thinking mode and reserve thinking for genuinely
hard problems}.
We formalize this intuition as \textsc{Town} (\textbf{T}hink
\textbf{O}nly \textbf{W}hen \textbf{N}eeded), a training-free,
two-stage inference cascade that requires no access to model internals
and no auxiliary models.

% ---------------------------------------------------------------------------
\subsection{Key Insight: Early Stopping as a Difficulty Classifier}
\label{sec:method:insight}
% ---------------------------------------------------------------------------

Our analysis reveals that non-thinking mode's early-stopping behavior
provides a natural binary classifier for problem difficulty:

\begin{observation}
\label{obs:early-stop}
For Qwen3-8B with \texttt{nothink@256} on GSM8K ($n{=}1{,}319$):
\begin{itemize}[nosep,leftmargin=*]
    \item \textbf{88.8\% of samples stop early} (avg 133 tokens),
      achieving \textbf{94.4\%} accuracy among early-stopped samples
      (87.5\% overall).
    \item \textbf{11.2\% of samples hit the budget} (use all 256 tokens),
      indicating the model struggled.
\end{itemize}
\end{observation}

This dichotomy provides a zero-cost routing signal: if the model finishes
quickly in non-thinking mode, the answer is likely correct; if it
exhausts the budget, the problem is likely hard and may benefit from
extended reasoning.
Unlike soft confidence scores or logit-based uncertainty measures, this
signal is \emph{binary}, \emph{endogenous}, and \emph{free}---it emerges
as a byproduct of the generation process itself
(Finding~2, \S\ref{sec:finding:natural-stop}).

% ---------------------------------------------------------------------------
\subsection{The \textsc{Town} Cascade}
\label{sec:method:cascade}
% ---------------------------------------------------------------------------

\textsc{Town} operates in two stages, summarized in
Algorithm~\ref{alg:town} and Figure~\ref{fig:town-pipeline}.

\begin{algorithm}[t]
\caption{\textsc{Town}: Think Only When Needed}
\label{alg:town}
\begin{algorithmic}[1]
\REQUIRE Question $q$, model $\mathcal{M}$,
  probe budget $B_1$ (non-thinking),
  fallback budget $B_2$ (thinking, $B_2 > B_1$)
\ENSURE Answer $\hat{a}$, total tokens consumed $T$

\medskip
\STATE \textbf{// Stage 1: Non-thinking probe}
\STATE $y_1 \leftarrow \mathcal{M}_{\textrm{nt}}(q,\; B_1)$
  \hfill\COMMENT{Generate without thinking}
\STATE $T \leftarrow |y_1|$

\medskip
\IF{$|y_1| < B_1$}
  \hfill\COMMENT{Natural stop $\Rightarrow$ high confidence}
  \STATE $\hat{a} \leftarrow a(y_1)$
  \hfill\COMMENT{Accept non-thinking answer}
\ELSE
  \hfill\COMMENT{Budget hit $\Rightarrow$ low confidence}
  \STATE \textbf{// Stage 2: Thinking fallback}
  \STATE $y_2 \leftarrow \mathcal{M}_{\textrm{th}}(q,\; B_2)$
    \hfill\COMMENT{Generate with chain-of-thought}
  \STATE $T \leftarrow T + |y_2|$
  \STATE $\hat{a} \leftarrow a(y_2)$
  \hfill\COMMENT{Use thinking answer}
\ENDIF

\medskip
\RETURN $\hat{a},\; T$
\end{algorithmic}
\end{algorithm}

% fig_town_pipeline.tex — TikZ pipeline diagram for TOWN method
% Include in method_final.tex via \input{sections/fig_town_pipeline}

\begin{figure}[t]
\centering
\resizebox{\linewidth}{!}{%
\begin{tikzpicture}[
    node distance=1.0cm and 1.2cm,
    box/.style={rectangle, draw, rounded corners=3pt, minimum width=2.0cm,
                minimum height=0.75cm, align=center, font=\small},
    decision/.style={diamond, draw, aspect=2.5, inner sep=1pt,
                     align=center, font=\small},
    arrow/.style={->, thick, >=stealth},
    label/.style={font=\scriptsize, text=gray},
]

% Input
\node[box, fill=gray!15] (input) {Input $q$};

% Stage 1
\node[box, fill=nothinkblue!20, right=1.5cm of input] (s1)
    {\textbf{Stage 1}\\Nothink@$B_1$};

% Decision
\node[decision, fill=yellow!15, right=1.5cm of s1] (dec)
    {$|y_1|<B_1$?};

% Accept
\node[box, fill=easycolor!25, above right=0.6cm and 1.5cm of dec] (accept)
    {\textbf{Accept}\\$\hat{a}=a(y_1)$};

% Stage 2
\node[box, fill=thinkorange!25, below right=0.6cm and 1.5cm of dec] (s2)
    {\textbf{Stage 2}\\Think@$B_2$};

% Final
\node[box, fill=thinkorange!15, right=1.2cm of s2] (final2)
    {\textbf{Accept}\\$\hat{a}=a(y_2)$};

% Arrows
\draw[arrow] (input) -- (s1);
\draw[arrow] (s1) -- (dec);
\draw[arrow] (dec) -- node[label, above left=-2pt] {Yes (88.8\%)} (accept);
\draw[arrow] (dec) -- node[label, below left=-2pt] {No (11.2\%)} (s2);
\draw[arrow] (s2) -- (final2);

% Annotations
\node[below=0.15cm of s1, font=\scriptsize, text=nothinkblue]
    {$\sim$133 tokens};
\node[above=0.15cm of accept, font=\scriptsize, text=easycolor!80!black]
    {94.4\% acc};
\node[below=0.15cm of s2, font=\scriptsize, text=thinkorange!80!black]
    {$\sim$469 tokens};
\node[above=0.15cm of final2, font=\scriptsize, text=thinkorange!80!black]
    {62.8\% acc};

\end{tikzpicture}%
}
\caption{%
    \textbf{\town{} baseline inference pipeline.}
    Stage~1 probes with non-thinking mode at budget $B_1$.
    If the model stops early (88.8\% of GSM8K), the answer is accepted
    at ${\sim}$133 tokens (94.4\% accuracy).
    Otherwise, Stage~2 routes to thinking mode at budget $B_2$,
    recovering additional correct answers on hard problems.
    Overall: \textbf{90.9\%} accuracy at \textbf{199} average tokens on the full test set ($n{=}1{,}319$).
    Note: \method{} extends this with decoupled answer generation and iterative refinement (Algorithm~\ref{alg:mrsd}).
}
\label{fig:town-pipeline}
\end{figure}


\paragraph{Stage 1: Non-thinking probe.}
We run the input through non-thinking mode with a moderate budget
$B_1$ (default: 256 tokens).
If the model produces a complete answer before exhausting $B_1$, we
accept it immediately.
This stage handles the vast majority of inputs---88.8\% on GSM8K---at
minimal cost (${\sim}$133 tokens on average for early-stopped samples), exploiting Finding~1's
observation that non-thinking mode is more token-efficient than thinking
mode at matched budgets.

\paragraph{Stage 2: Thinking fallback.}
For the ${\sim}$11.2\% of inputs that exhaust the Stage~1 budget---indicating
the model could not produce a confident answer quickly---we route to
thinking mode with an extended budget $B_2$ (default: 512 tokens).
These are precisely the ``hard'' problems where step-by-step reasoning
chains are most likely to provide genuine benefit.
The total cost for a fallback sample is $B_1 + |y_2|$, since the
probe consumed its full budget.

% ---------------------------------------------------------------------------
\subsection{Cost and Accuracy Analysis}
\label{sec:method:analysis}
% ---------------------------------------------------------------------------

\paragraph{Expected token cost.}
Let $p$ denote the early-stop rate in Stage~1,
$\bar{t}_1$ the average tokens for early-stopped samples, and
$\bar{t}_2$ the average thinking-mode cost for fallback samples.
The expected token cost per question is:
\begin{equation}
  \label{eq:town-cost}
  \mathbb{E}[T] \;=\; p \cdot \bar{t}_1
    \;+\; (1 - p) \cdot \big(B_1 + \bar{t}_2\big).
\end{equation}
With our GSM8K parameters ($p{=}0.888$, $\bar{t}_1{\approx}133$,
$B_1{=}256$, $\bar{t}_2{\approx}469$) on the full test set ($n{=}1{,}319$):
\begin{equation}
  \label{eq:town-cost-concrete}
  \mathbb{E}[T]
    \;\approx\; 0.888 \times 133 \;+\; 0.112 \times (256 + 469)
    \;\approx\; 118 + 81
    \;=\; \mathbf{\sim\!199\text{ tokens}}.
\end{equation}
This is \textbf{2.4$\times$} cheaper than \texttt{think@512} alone
(460~tokens), representing a \textbf{57\%} reduction in average token
consumption.  Note that the Stage~1 probe budget $B_1$ is ``spent''
even for routed samples---it serves as the cost of difficulty
classification.

\paragraph{Projected accuracy.}
\textsc{Town}'s accuracy decomposes as:
\begin{equation}
  \label{eq:town-acc}
  \textrm{Acc}_{\textsc{Town}}
    \;=\; p \cdot \textrm{Acc}_{\textrm{nt}}^{\textrm{early}}
    \;+\; (1 - p) \cdot \textrm{Acc}_{\textrm{th}}^{\textrm{hard}},
\end{equation}
where $\textrm{Acc}_{\textrm{nt}}^{\textrm{early}}$ is the
accuracy of early-stop non-thinking samples and
$\textrm{Acc}_{\textrm{th}}^{\textrm{hard}}$ is the thinking-mode
accuracy on the hard subset routed to Stage~2.
On the full GSM8K test set ($n{=}1{,}319$),
\textsc{Town} achieves \textbf{90.9\%} accuracy at only \textbf{199
average tokens} (including the Stage~1 probe cost for routed samples):
\begin{equation}
  \textrm{Acc}_{\textsc{Town}}
    \;=\; \underbrace{0.888 \times 94.4\%}_{\text{Stage 1: early-stop}}
    \;+\; \underbrace{0.112 \times 62.8\%}_{\text{Stage 2: thinking fallback}}
    \;=\; 83.8\% + 7.0\%
    \;=\; \mathbf{90.9\%}.
\end{equation}
This exceeds \texttt{think@512} (65.2\%) by \textbf{+25.7\,pp}
while using \textbf{2.4$\times$ fewer tokens}, and surpasses
\texttt{nothink@256} (87.5\%) by \textbf{+3.4\,pp} by recovering
hard cases through Stage~2.
Full evaluation results are in \S\ref{sec:town-eval}.

% ---------------------------------------------------------------------------
\subsection{Comparison with Alternative Strategies}
\label{sec:method:comparison}
% ---------------------------------------------------------------------------

\paragraph{vs.\ \texttt{think@512} only.}
\textsc{Town} achieves substantially higher accuracy at
${\sim}$2.4$\times$ fewer tokens.
The advantage comes from avoiding thinking mode's truncation waste on the
88.8\% of problems where it is unnecessary.

\paragraph{vs.\ \texttt{nothink@256} only.}
\textsc{Town} adds a thinking fallback for the $\sim$11.2\% of hard
problems, recovering accuracy on cases where direct answering
fails.
The additional cost is modest: $\sim$53 extra tokens on average
over the entire dataset (199 vs.\ 146 for \texttt{nothink@256}).

\paragraph{vs.\ Self-consistency~\cite{wang2023selfconsistency}.}
SC@$K$ generates $K$ full-budget reasoning traces and majority-votes,
consuming $K \times B$ tokens unconditionally.
\textsc{Town} uses a \emph{single} non-thinking probe and a binary
routing decision, yielding lower latency (one generation call for easy
questions vs.\ $K$), lower token cost ($\sim$146 tokens for easy
questions vs.\ $K \times B$), and a simpler implementation (no consensus
computation, no threshold tuning).

\paragraph{vs.\ Oracle routing.}
An oracle that knows each problem's difficulty \emph{a priori}
achieves 68.2\% accuracy at 401 average tokens
(Table~\ref{tab:efficiency-frontier}).
\textsc{Town}'s early-stop signal approximates this oracle using a
training-free heuristic, with the gap determined primarily by the
signal's precision on the routed subset.

\paragraph{Hyperparameters.}
\textsc{Town} introduces exactly \emph{two} hyperparameters beyond the
base model: $B_1$ (probe budget) and $B_2$ (thinking fallback budget).
Both have intuitive interpretations: $B_1$ should be large enough for
the model to produce a complete non-thinking answer on most easy
questions (256 tokens suffices for GSM8K), and $B_2$ should be large
enough for a complete reasoning chain on hard questions (512+ for
GSM8K).
No routing thresholds, consensus weights, or learned parameters are
required.

% ===========================================================================

The empirical findings in \S\ref{sec:thinking-tax} suggest a clear strategy:
\textbf{default to non-thinking mode and reserve thinking for genuinely
hard problems}.
We formalize this intuition as \textsc{Town} (\textbf{T}hink
\textbf{O}nly \textbf{W}hen \textbf{N}eeded), a training-free,
two-stage inference cascade that requires no access to model internals
and no auxiliary models.

% ---------------------------------------------------------------------------
\subsection{Key Insight: Early Stopping as a Difficulty Classifier}
\label{sec:method:insight}
% ---------------------------------------------------------------------------

Our analysis reveals that non-thinking mode's early-stopping behavior
provides a natural binary classifier for problem difficulty:

\begin{observation}
\label{obs:early-stop}
For Qwen3-8B with \texttt{nothink@256} on GSM8K ($n{=}1{,}319$):
\begin{itemize}[nosep,leftmargin=*]
    \item \textbf{88.8\% of samples stop early} (avg 133 tokens),
      achieving \textbf{94.4\%} accuracy among early-stopped samples
      (87.5\% overall).
    \item \textbf{11.2\% of samples hit the budget} (use all 256 tokens),
      indicating the model struggled.
\end{itemize}
\end{observation}

This dichotomy provides a zero-cost routing signal: if the model finishes
quickly in non-thinking mode, the answer is likely correct; if it
exhausts the budget, the problem is likely hard and may benefit from
extended reasoning.
Unlike soft confidence scores or logit-based uncertainty measures, this
signal is \emph{binary}, \emph{endogenous}, and \emph{free}---it emerges
as a byproduct of the generation process itself
(Finding~2, \S\ref{sec:finding:natural-stop}).

% ---------------------------------------------------------------------------
\subsection{The \textsc{Town} Cascade}
\label{sec:method:cascade}
% ---------------------------------------------------------------------------

\textsc{Town} operates in two stages, summarized in
Algorithm~\ref{alg:town} and Figure~\ref{fig:town-pipeline}.

\begin{algorithm}[t]
\caption{\textsc{Town}: Think Only When Needed}
\label{alg:town}
\begin{algorithmic}[1]
\REQUIRE Question $q$, model $\mathcal{M}$,
  probe budget $B_1$ (non-thinking),
  fallback budget $B_2$ (thinking, $B_2 > B_1$)
\ENSURE Answer $\hat{a}$, total tokens consumed $T$

\medskip
\STATE \textbf{// Stage 1: Non-thinking probe}
\STATE $y_1 \leftarrow \mathcal{M}_{\textrm{nt}}(q,\; B_1)$
  \hfill\COMMENT{Generate without thinking}
\STATE $T \leftarrow |y_1|$

\medskip
\IF{$|y_1| < B_1$}
  \hfill\COMMENT{Natural stop $\Rightarrow$ high confidence}
  \STATE $\hat{a} \leftarrow a(y_1)$
  \hfill\COMMENT{Accept non-thinking answer}
\ELSE
  \hfill\COMMENT{Budget hit $\Rightarrow$ low confidence}
  \STATE \textbf{// Stage 2: Thinking fallback}
  \STATE $y_2 \leftarrow \mathcal{M}_{\textrm{th}}(q,\; B_2)$
    \hfill\COMMENT{Generate with chain-of-thought}
  \STATE $T \leftarrow T + |y_2|$
  \STATE $\hat{a} \leftarrow a(y_2)$
  \hfill\COMMENT{Use thinking answer}
\ENDIF

\medskip
\RETURN $\hat{a},\; T$
\end{algorithmic}
\end{algorithm}

% fig_town_pipeline.tex — TikZ pipeline diagram for TOWN method
% Include in method_final.tex via % fig_town_pipeline.tex — TikZ pipeline diagram for TOWN method
% Include in method_final.tex via \input{sections/fig_town_pipeline}

\begin{figure}[t]
\centering
\resizebox{\linewidth}{!}{%
\begin{tikzpicture}[
    node distance=1.0cm and 1.2cm,
    box/.style={rectangle, draw, rounded corners=3pt, minimum width=2.0cm,
                minimum height=0.75cm, align=center, font=\small},
    decision/.style={diamond, draw, aspect=2.5, inner sep=1pt,
                     align=center, font=\small},
    arrow/.style={->, thick, >=stealth},
    label/.style={font=\scriptsize, text=gray},
]

% Input
\node[box, fill=gray!15] (input) {Input $q$};

% Stage 1
\node[box, fill=nothinkblue!20, right=1.5cm of input] (s1)
    {\textbf{Stage 1}\\Nothink@$B_1$};

% Decision
\node[decision, fill=yellow!15, right=1.5cm of s1] (dec)
    {$|y_1|<B_1$?};

% Accept
\node[box, fill=easycolor!25, above right=0.6cm and 1.5cm of dec] (accept)
    {\textbf{Accept}\\$\hat{a}=a(y_1)$};

% Stage 2
\node[box, fill=thinkorange!25, below right=0.6cm and 1.5cm of dec] (s2)
    {\textbf{Stage 2}\\Think@$B_2$};

% Final
\node[box, fill=thinkorange!15, right=1.2cm of s2] (final2)
    {\textbf{Accept}\\$\hat{a}=a(y_2)$};

% Arrows
\draw[arrow] (input) -- (s1);
\draw[arrow] (s1) -- (dec);
\draw[arrow] (dec) -- node[label, above left=-2pt] {Yes (88.8\%)} (accept);
\draw[arrow] (dec) -- node[label, below left=-2pt] {No (11.2\%)} (s2);
\draw[arrow] (s2) -- (final2);

% Annotations
\node[below=0.15cm of s1, font=\scriptsize, text=nothinkblue]
    {$\sim$133 tokens};
\node[above=0.15cm of accept, font=\scriptsize, text=easycolor!80!black]
    {94.4\% acc};
\node[below=0.15cm of s2, font=\scriptsize, text=thinkorange!80!black]
    {$\sim$469 tokens};
\node[above=0.15cm of final2, font=\scriptsize, text=thinkorange!80!black]
    {62.8\% acc};

\end{tikzpicture}%
}
\caption{%
    \textbf{\town{} baseline inference pipeline.}
    Stage~1 probes with non-thinking mode at budget $B_1$.
    If the model stops early (88.8\% of GSM8K), the answer is accepted
    at ${\sim}$133 tokens (94.4\% accuracy).
    Otherwise, Stage~2 routes to thinking mode at budget $B_2$,
    recovering additional correct answers on hard problems.
    Overall: \textbf{90.9\%} accuracy at \textbf{199} average tokens on the full test set ($n{=}1{,}319$).
    Note: \method{} extends this with decoupled answer generation and iterative refinement (Algorithm~\ref{alg:mrsd}).
}
\label{fig:town-pipeline}
\end{figure}


\begin{figure}[t]
\centering
\resizebox{\linewidth}{!}{%
\begin{tikzpicture}[
    node distance=1.0cm and 1.2cm,
    box/.style={rectangle, draw, rounded corners=3pt, minimum width=2.0cm,
                minimum height=0.75cm, align=center, font=\small},
    decision/.style={diamond, draw, aspect=2.5, inner sep=1pt,
                     align=center, font=\small},
    arrow/.style={->, thick, >=stealth},
    label/.style={font=\scriptsize, text=gray},
]

% Input
\node[box, fill=gray!15] (input) {Input $q$};

% Stage 1
\node[box, fill=nothinkblue!20, right=1.5cm of input] (s1)
    {\textbf{Stage 1}\\Nothink@$B_1$};

% Decision
\node[decision, fill=yellow!15, right=1.5cm of s1] (dec)
    {$|y_1|<B_1$?};

% Accept
\node[box, fill=easycolor!25, above right=0.6cm and 1.5cm of dec] (accept)
    {\textbf{Accept}\\$\hat{a}=a(y_1)$};

% Stage 2
\node[box, fill=thinkorange!25, below right=0.6cm and 1.5cm of dec] (s2)
    {\textbf{Stage 2}\\Think@$B_2$};

% Final
\node[box, fill=thinkorange!15, right=1.2cm of s2] (final2)
    {\textbf{Accept}\\$\hat{a}=a(y_2)$};

% Arrows
\draw[arrow] (input) -- (s1);
\draw[arrow] (s1) -- (dec);
\draw[arrow] (dec) -- node[label, above left=-2pt] {Yes (88.8\%)} (accept);
\draw[arrow] (dec) -- node[label, below left=-2pt] {No (11.2\%)} (s2);
\draw[arrow] (s2) -- (final2);

% Annotations
\node[below=0.15cm of s1, font=\scriptsize, text=nothinkblue]
    {$\sim$133 tokens};
\node[above=0.15cm of accept, font=\scriptsize, text=easycolor!80!black]
    {94.4\% acc};
\node[below=0.15cm of s2, font=\scriptsize, text=thinkorange!80!black]
    {$\sim$469 tokens};
\node[above=0.15cm of final2, font=\scriptsize, text=thinkorange!80!black]
    {62.8\% acc};

\end{tikzpicture}%
}
\caption{%
    \textbf{\town{} baseline inference pipeline.}
    Stage~1 probes with non-thinking mode at budget $B_1$.
    If the model stops early (88.8\% of GSM8K), the answer is accepted
    at ${\sim}$133 tokens (94.4\% accuracy).
    Otherwise, Stage~2 routes to thinking mode at budget $B_2$,
    recovering additional correct answers on hard problems.
    Overall: \textbf{90.9\%} accuracy at \textbf{199} average tokens on the full test set ($n{=}1{,}319$).
    Note: \method{} extends this with decoupled answer generation and iterative refinement (Algorithm~\ref{alg:mrsd}).
}
\label{fig:town-pipeline}
\end{figure}


\paragraph{Stage 1: Non-thinking probe.}
We run the input through non-thinking mode with a moderate budget
$B_1$ (default: 256 tokens).
If the model produces a complete answer before exhausting $B_1$, we
accept it immediately.
This stage handles the vast majority of inputs---88.8\% on GSM8K---at
minimal cost (${\sim}$133 tokens on average for early-stopped samples), exploiting Finding~1's
observation that non-thinking mode is more token-efficient than thinking
mode at matched budgets.

\paragraph{Stage 2: Thinking fallback.}
For the ${\sim}$11.2\% of inputs that exhaust the Stage~1 budget---indicating
the model could not produce a confident answer quickly---we route to
thinking mode with an extended budget $B_2$ (default: 512 tokens).
These are precisely the ``hard'' problems where step-by-step reasoning
chains are most likely to provide genuine benefit.
The total cost for a fallback sample is $B_1 + |y_2|$, since the
probe consumed its full budget.

% ---------------------------------------------------------------------------
\subsection{Cost and Accuracy Analysis}
\label{sec:method:analysis}
% ---------------------------------------------------------------------------

\paragraph{Expected token cost.}
Let $p$ denote the early-stop rate in Stage~1,
$\bar{t}_1$ the average tokens for early-stopped samples, and
$\bar{t}_2$ the average thinking-mode cost for fallback samples.
The expected token cost per question is:
\begin{equation}
  \label{eq:town-cost}
  \mathbb{E}[T] \;=\; p \cdot \bar{t}_1
    \;+\; (1 - p) \cdot \big(B_1 + \bar{t}_2\big).
\end{equation}
With our GSM8K parameters ($p{=}0.888$, $\bar{t}_1{\approx}133$,
$B_1{=}256$, $\bar{t}_2{\approx}469$) on the full test set ($n{=}1{,}319$):
\begin{equation}
  \label{eq:town-cost-concrete}
  \mathbb{E}[T]
    \;\approx\; 0.888 \times 133 \;+\; 0.112 \times (256 + 469)
    \;\approx\; 118 + 81
    \;=\; \mathbf{\sim\!199\text{ tokens}}.
\end{equation}
This is \textbf{2.4$\times$} cheaper than \texttt{think@512} alone
(460~tokens), representing a \textbf{57\%} reduction in average token
consumption.  Note that the Stage~1 probe budget $B_1$ is ``spent''
even for routed samples---it serves as the cost of difficulty
classification.

\paragraph{Projected accuracy.}
\textsc{Town}'s accuracy decomposes as:
\begin{equation}
  \label{eq:town-acc}
  \textrm{Acc}_{\textsc{Town}}
    \;=\; p \cdot \textrm{Acc}_{\textrm{nt}}^{\textrm{early}}
    \;+\; (1 - p) \cdot \textrm{Acc}_{\textrm{th}}^{\textrm{hard}},
\end{equation}
where $\textrm{Acc}_{\textrm{nt}}^{\textrm{early}}$ is the
accuracy of early-stop non-thinking samples and
$\textrm{Acc}_{\textrm{th}}^{\textrm{hard}}$ is the thinking-mode
accuracy on the hard subset routed to Stage~2.
On the full GSM8K test set ($n{=}1{,}319$),
\textsc{Town} achieves \textbf{90.9\%} accuracy at only \textbf{199
average tokens} (including the Stage~1 probe cost for routed samples):
\begin{equation}
  \textrm{Acc}_{\textsc{Town}}
    \;=\; \underbrace{0.888 \times 94.4\%}_{\text{Stage 1: early-stop}}
    \;+\; \underbrace{0.112 \times 62.8\%}_{\text{Stage 2: thinking fallback}}
    \;=\; 83.8\% + 7.0\%
    \;=\; \mathbf{90.9\%}.
\end{equation}
This exceeds \texttt{think@512} (65.2\%) by \textbf{+25.7\,pp}
while using \textbf{2.4$\times$ fewer tokens}, and surpasses
\texttt{nothink@256} (87.5\%) by \textbf{+3.4\,pp} by recovering
hard cases through Stage~2.
Full evaluation results are in \S\ref{sec:town-eval}.

% ---------------------------------------------------------------------------
\subsection{Comparison with Alternative Strategies}
\label{sec:method:comparison}
% ---------------------------------------------------------------------------

\paragraph{vs.\ \texttt{think@512} only.}
\textsc{Town} achieves substantially higher accuracy at
${\sim}$2.4$\times$ fewer tokens.
The advantage comes from avoiding thinking mode's truncation waste on the
88.8\% of problems where it is unnecessary.

\paragraph{vs.\ \texttt{nothink@256} only.}
\textsc{Town} adds a thinking fallback for the $\sim$11.2\% of hard
problems, recovering accuracy on cases where direct answering
fails.
The additional cost is modest: $\sim$53 extra tokens on average
over the entire dataset (199 vs.\ 146 for \texttt{nothink@256}).

\paragraph{vs.\ Self-consistency~\cite{wang2023selfconsistency}.}
SC@$K$ generates $K$ full-budget reasoning traces and majority-votes,
consuming $K \times B$ tokens unconditionally.
\textsc{Town} uses a \emph{single} non-thinking probe and a binary
routing decision, yielding lower latency (one generation call for easy
questions vs.\ $K$), lower token cost ($\sim$146 tokens for easy
questions vs.\ $K \times B$), and a simpler implementation (no consensus
computation, no threshold tuning).

\paragraph{vs.\ Oracle routing.}
An oracle that knows each problem's difficulty \emph{a priori}
achieves 68.2\% accuracy at 401 average tokens
(Table~\ref{tab:efficiency-frontier}).
\textsc{Town}'s early-stop signal approximates this oracle using a
training-free heuristic, with the gap determined primarily by the
signal's precision on the routed subset.

\paragraph{Hyperparameters.}
\textsc{Town} introduces exactly \emph{two} hyperparameters beyond the
base model: $B_1$ (probe budget) and $B_2$ (thinking fallback budget).
Both have intuitive interpretations: $B_1$ should be large enough for
the model to produce a complete non-thinking answer on most easy
questions (256 tokens suffices for GSM8K), and $B_2$ should be large
enough for a complete reasoning chain on hard questions (512+ for
GSM8K).
No routing thresholds, consensus weights, or learned parameters are
required.

% ===========================================================================

The empirical findings in \S\ref{sec:thinking-tax} suggest a clear strategy:
\textbf{default to non-thinking mode and reserve thinking for genuinely
hard problems}.
We formalize this intuition as \textsc{Town} (\textbf{T}hink
\textbf{O}nly \textbf{W}hen \textbf{N}eeded), a training-free,
two-stage inference cascade that requires no access to model internals
and no auxiliary models.

% ---------------------------------------------------------------------------
\subsection{Key Insight: Early Stopping as a Difficulty Classifier}
\label{sec:method:insight}
% ---------------------------------------------------------------------------

Our analysis reveals that non-thinking mode's early-stopping behavior
provides a natural binary classifier for problem difficulty:

\begin{observation}
\label{obs:early-stop}
For Qwen3-8B with \texttt{nothink@256} on GSM8K ($n{=}1{,}319$):
\begin{itemize}[nosep,leftmargin=*]
    \item \textbf{88.8\% of samples stop early} (avg 133 tokens),
      achieving \textbf{94.4\%} accuracy among early-stopped samples
      (87.5\% overall).
    \item \textbf{11.2\% of samples hit the budget} (use all 256 tokens),
      indicating the model struggled.
\end{itemize}
\end{observation}

This dichotomy provides a zero-cost routing signal: if the model finishes
quickly in non-thinking mode, the answer is likely correct; if it
exhausts the budget, the problem is likely hard and may benefit from
extended reasoning.
Unlike soft confidence scores or logit-based uncertainty measures, this
signal is \emph{binary}, \emph{endogenous}, and \emph{free}---it emerges
as a byproduct of the generation process itself
(Finding~2, \S\ref{sec:finding:natural-stop}).

% ---------------------------------------------------------------------------
\subsection{The \textsc{Town} Cascade}
\label{sec:method:cascade}
% ---------------------------------------------------------------------------

\textsc{Town} operates in two stages, summarized in
Algorithm~\ref{alg:town} and Figure~\ref{fig:town-pipeline}.

\begin{algorithm}[t]
\caption{\textsc{Town}: Think Only When Needed}
\label{alg:town}
\begin{algorithmic}[1]
\REQUIRE Question $q$, model $\mathcal{M}$,
  probe budget $B_1$ (non-thinking),
  fallback budget $B_2$ (thinking, $B_2 > B_1$)
\ENSURE Answer $\hat{a}$, total tokens consumed $T$

\medskip
\STATE \textbf{// Stage 1: Non-thinking probe}
\STATE $y_1 \leftarrow \mathcal{M}_{\textrm{nt}}(q,\; B_1)$
  \hfill\COMMENT{Generate without thinking}
\STATE $T \leftarrow |y_1|$

\medskip
\IF{$|y_1| < B_1$}
  \hfill\COMMENT{Natural stop $\Rightarrow$ high confidence}
  \STATE $\hat{a} \leftarrow a(y_1)$
  \hfill\COMMENT{Accept non-thinking answer}
\ELSE
  \hfill\COMMENT{Budget hit $\Rightarrow$ low confidence}
  \STATE \textbf{// Stage 2: Thinking fallback}
  \STATE $y_2 \leftarrow \mathcal{M}_{\textrm{th}}(q,\; B_2)$
    \hfill\COMMENT{Generate with chain-of-thought}
  \STATE $T \leftarrow T + |y_2|$
  \STATE $\hat{a} \leftarrow a(y_2)$
  \hfill\COMMENT{Use thinking answer}
\ENDIF

\medskip
\RETURN $\hat{a},\; T$
\end{algorithmic}
\end{algorithm}

% fig_town_pipeline.tex — TikZ pipeline diagram for TOWN method
% Include in method_final.tex via % fig_town_pipeline.tex — TikZ pipeline diagram for TOWN method
% Include in method_final.tex via % fig_town_pipeline.tex — TikZ pipeline diagram for TOWN method
% Include in method_final.tex via \input{sections/fig_town_pipeline}

\begin{figure}[t]
\centering
\resizebox{\linewidth}{!}{%
\begin{tikzpicture}[
    node distance=1.0cm and 1.2cm,
    box/.style={rectangle, draw, rounded corners=3pt, minimum width=2.0cm,
                minimum height=0.75cm, align=center, font=\small},
    decision/.style={diamond, draw, aspect=2.5, inner sep=1pt,
                     align=center, font=\small},
    arrow/.style={->, thick, >=stealth},
    label/.style={font=\scriptsize, text=gray},
]

% Input
\node[box, fill=gray!15] (input) {Input $q$};

% Stage 1
\node[box, fill=nothinkblue!20, right=1.5cm of input] (s1)
    {\textbf{Stage 1}\\Nothink@$B_1$};

% Decision
\node[decision, fill=yellow!15, right=1.5cm of s1] (dec)
    {$|y_1|<B_1$?};

% Accept
\node[box, fill=easycolor!25, above right=0.6cm and 1.5cm of dec] (accept)
    {\textbf{Accept}\\$\hat{a}=a(y_1)$};

% Stage 2
\node[box, fill=thinkorange!25, below right=0.6cm and 1.5cm of dec] (s2)
    {\textbf{Stage 2}\\Think@$B_2$};

% Final
\node[box, fill=thinkorange!15, right=1.2cm of s2] (final2)
    {\textbf{Accept}\\$\hat{a}=a(y_2)$};

% Arrows
\draw[arrow] (input) -- (s1);
\draw[arrow] (s1) -- (dec);
\draw[arrow] (dec) -- node[label, above left=-2pt] {Yes (88.8\%)} (accept);
\draw[arrow] (dec) -- node[label, below left=-2pt] {No (11.2\%)} (s2);
\draw[arrow] (s2) -- (final2);

% Annotations
\node[below=0.15cm of s1, font=\scriptsize, text=nothinkblue]
    {$\sim$133 tokens};
\node[above=0.15cm of accept, font=\scriptsize, text=easycolor!80!black]
    {94.4\% acc};
\node[below=0.15cm of s2, font=\scriptsize, text=thinkorange!80!black]
    {$\sim$469 tokens};
\node[above=0.15cm of final2, font=\scriptsize, text=thinkorange!80!black]
    {62.8\% acc};

\end{tikzpicture}%
}
\caption{%
    \textbf{\town{} baseline inference pipeline.}
    Stage~1 probes with non-thinking mode at budget $B_1$.
    If the model stops early (88.8\% of GSM8K), the answer is accepted
    at ${\sim}$133 tokens (94.4\% accuracy).
    Otherwise, Stage~2 routes to thinking mode at budget $B_2$,
    recovering additional correct answers on hard problems.
    Overall: \textbf{90.9\%} accuracy at \textbf{199} average tokens on the full test set ($n{=}1{,}319$).
    Note: \method{} extends this with decoupled answer generation and iterative refinement (Algorithm~\ref{alg:mrsd}).
}
\label{fig:town-pipeline}
\end{figure}


\begin{figure}[t]
\centering
\resizebox{\linewidth}{!}{%
\begin{tikzpicture}[
    node distance=1.0cm and 1.2cm,
    box/.style={rectangle, draw, rounded corners=3pt, minimum width=2.0cm,
                minimum height=0.75cm, align=center, font=\small},
    decision/.style={diamond, draw, aspect=2.5, inner sep=1pt,
                     align=center, font=\small},
    arrow/.style={->, thick, >=stealth},
    label/.style={font=\scriptsize, text=gray},
]

% Input
\node[box, fill=gray!15] (input) {Input $q$};

% Stage 1
\node[box, fill=nothinkblue!20, right=1.5cm of input] (s1)
    {\textbf{Stage 1}\\Nothink@$B_1$};

% Decision
\node[decision, fill=yellow!15, right=1.5cm of s1] (dec)
    {$|y_1|<B_1$?};

% Accept
\node[box, fill=easycolor!25, above right=0.6cm and 1.5cm of dec] (accept)
    {\textbf{Accept}\\$\hat{a}=a(y_1)$};

% Stage 2
\node[box, fill=thinkorange!25, below right=0.6cm and 1.5cm of dec] (s2)
    {\textbf{Stage 2}\\Think@$B_2$};

% Final
\node[box, fill=thinkorange!15, right=1.2cm of s2] (final2)
    {\textbf{Accept}\\$\hat{a}=a(y_2)$};

% Arrows
\draw[arrow] (input) -- (s1);
\draw[arrow] (s1) -- (dec);
\draw[arrow] (dec) -- node[label, above left=-2pt] {Yes (88.8\%)} (accept);
\draw[arrow] (dec) -- node[label, below left=-2pt] {No (11.2\%)} (s2);
\draw[arrow] (s2) -- (final2);

% Annotations
\node[below=0.15cm of s1, font=\scriptsize, text=nothinkblue]
    {$\sim$133 tokens};
\node[above=0.15cm of accept, font=\scriptsize, text=easycolor!80!black]
    {94.4\% acc};
\node[below=0.15cm of s2, font=\scriptsize, text=thinkorange!80!black]
    {$\sim$469 tokens};
\node[above=0.15cm of final2, font=\scriptsize, text=thinkorange!80!black]
    {62.8\% acc};

\end{tikzpicture}%
}
\caption{%
    \textbf{\town{} baseline inference pipeline.}
    Stage~1 probes with non-thinking mode at budget $B_1$.
    If the model stops early (88.8\% of GSM8K), the answer is accepted
    at ${\sim}$133 tokens (94.4\% accuracy).
    Otherwise, Stage~2 routes to thinking mode at budget $B_2$,
    recovering additional correct answers on hard problems.
    Overall: \textbf{90.9\%} accuracy at \textbf{199} average tokens on the full test set ($n{=}1{,}319$).
    Note: \method{} extends this with decoupled answer generation and iterative refinement (Algorithm~\ref{alg:mrsd}).
}
\label{fig:town-pipeline}
\end{figure}


\begin{figure}[t]
\centering
\resizebox{\linewidth}{!}{%
\begin{tikzpicture}[
    node distance=1.0cm and 1.2cm,
    box/.style={rectangle, draw, rounded corners=3pt, minimum width=2.0cm,
                minimum height=0.75cm, align=center, font=\small},
    decision/.style={diamond, draw, aspect=2.5, inner sep=1pt,
                     align=center, font=\small},
    arrow/.style={->, thick, >=stealth},
    label/.style={font=\scriptsize, text=gray},
]

% Input
\node[box, fill=gray!15] (input) {Input $q$};

% Stage 1
\node[box, fill=nothinkblue!20, right=1.5cm of input] (s1)
    {\textbf{Stage 1}\\Nothink@$B_1$};

% Decision
\node[decision, fill=yellow!15, right=1.5cm of s1] (dec)
    {$|y_1|<B_1$?};

% Accept
\node[box, fill=easycolor!25, above right=0.6cm and 1.5cm of dec] (accept)
    {\textbf{Accept}\\$\hat{a}=a(y_1)$};

% Stage 2
\node[box, fill=thinkorange!25, below right=0.6cm and 1.5cm of dec] (s2)
    {\textbf{Stage 2}\\Think@$B_2$};

% Final
\node[box, fill=thinkorange!15, right=1.2cm of s2] (final2)
    {\textbf{Accept}\\$\hat{a}=a(y_2)$};

% Arrows
\draw[arrow] (input) -- (s1);
\draw[arrow] (s1) -- (dec);
\draw[arrow] (dec) -- node[label, above left=-2pt] {Yes (88.8\%)} (accept);
\draw[arrow] (dec) -- node[label, below left=-2pt] {No (11.2\%)} (s2);
\draw[arrow] (s2) -- (final2);

% Annotations
\node[below=0.15cm of s1, font=\scriptsize, text=nothinkblue]
    {$\sim$133 tokens};
\node[above=0.15cm of accept, font=\scriptsize, text=easycolor!80!black]
    {94.4\% acc};
\node[below=0.15cm of s2, font=\scriptsize, text=thinkorange!80!black]
    {$\sim$469 tokens};
\node[above=0.15cm of final2, font=\scriptsize, text=thinkorange!80!black]
    {62.8\% acc};

\end{tikzpicture}%
}
\caption{%
    \textbf{\town{} baseline inference pipeline.}
    Stage~1 probes with non-thinking mode at budget $B_1$.
    If the model stops early (88.8\% of GSM8K), the answer is accepted
    at ${\sim}$133 tokens (94.4\% accuracy).
    Otherwise, Stage~2 routes to thinking mode at budget $B_2$,
    recovering additional correct answers on hard problems.
    Overall: \textbf{90.9\%} accuracy at \textbf{199} average tokens on the full test set ($n{=}1{,}319$).
    Note: \method{} extends this with decoupled answer generation and iterative refinement (Algorithm~\ref{alg:mrsd}).
}
\label{fig:town-pipeline}
\end{figure}


\paragraph{Stage 1: Non-thinking probe.}
We run the input through non-thinking mode with a moderate budget
$B_1$ (default: 256 tokens).
If the model produces a complete answer before exhausting $B_1$, we
accept it immediately.
This stage handles the vast majority of inputs---88.8\% on GSM8K---at
minimal cost (${\sim}$133 tokens on average for early-stopped samples), exploiting Finding~1's
observation that non-thinking mode is more token-efficient than thinking
mode at matched budgets.

\paragraph{Stage 2: Thinking fallback.}
For the ${\sim}$11.2\% of inputs that exhaust the Stage~1 budget---indicating
the model could not produce a confident answer quickly---we route to
thinking mode with an extended budget $B_2$ (default: 512 tokens).
These are precisely the ``hard'' problems where step-by-step reasoning
chains are most likely to provide genuine benefit.
The total cost for a fallback sample is $B_1 + |y_2|$, since the
probe consumed its full budget.

% ---------------------------------------------------------------------------
\subsection{Cost and Accuracy Analysis}
\label{sec:method:analysis}
% ---------------------------------------------------------------------------

\paragraph{Expected token cost.}
Let $p$ denote the early-stop rate in Stage~1,
$\bar{t}_1$ the average tokens for early-stopped samples, and
$\bar{t}_2$ the average thinking-mode cost for fallback samples.
The expected token cost per question is:
\begin{equation}
  \label{eq:town-cost}
  \mathbb{E}[T] \;=\; p \cdot \bar{t}_1
    \;+\; (1 - p) \cdot \big(B_1 + \bar{t}_2\big).
\end{equation}
With our GSM8K parameters ($p{=}0.888$, $\bar{t}_1{\approx}133$,
$B_1{=}256$, $\bar{t}_2{\approx}469$) on the full test set ($n{=}1{,}319$):
\begin{equation}
  \label{eq:town-cost-concrete}
  \mathbb{E}[T]
    \;\approx\; 0.888 \times 133 \;+\; 0.112 \times (256 + 469)
    \;\approx\; 118 + 81
    \;=\; \mathbf{\sim\!199\text{ tokens}}.
\end{equation}
This is \textbf{2.4$\times$} cheaper than \texttt{think@512} alone
(460~tokens), representing a \textbf{57\%} reduction in average token
consumption.  Note that the Stage~1 probe budget $B_1$ is ``spent''
even for routed samples---it serves as the cost of difficulty
classification.

\paragraph{Projected accuracy.}
\textsc{Town}'s accuracy decomposes as:
\begin{equation}
  \label{eq:town-acc}
  \textrm{Acc}_{\textsc{Town}}
    \;=\; p \cdot \textrm{Acc}_{\textrm{nt}}^{\textrm{early}}
    \;+\; (1 - p) \cdot \textrm{Acc}_{\textrm{th}}^{\textrm{hard}},
\end{equation}
where $\textrm{Acc}_{\textrm{nt}}^{\textrm{early}}$ is the
accuracy of early-stop non-thinking samples and
$\textrm{Acc}_{\textrm{th}}^{\textrm{hard}}$ is the thinking-mode
accuracy on the hard subset routed to Stage~2.
On the full GSM8K test set ($n{=}1{,}319$),
\textsc{Town} achieves \textbf{90.9\%} accuracy at only \textbf{199
average tokens} (including the Stage~1 probe cost for routed samples):
\begin{equation}
  \textrm{Acc}_{\textsc{Town}}
    \;=\; \underbrace{0.888 \times 94.4\%}_{\text{Stage 1: early-stop}}
    \;+\; \underbrace{0.112 \times 62.8\%}_{\text{Stage 2: thinking fallback}}
    \;=\; 83.8\% + 7.0\%
    \;=\; \mathbf{90.9\%}.
\end{equation}
This exceeds \texttt{think@512} (65.2\%) by \textbf{+25.7\,pp}
while using \textbf{2.4$\times$ fewer tokens}, and surpasses
\texttt{nothink@256} (87.5\%) by \textbf{+3.4\,pp} by recovering
hard cases through Stage~2.
Full evaluation results are in \S\ref{sec:town-eval}.

% ---------------------------------------------------------------------------
\subsection{Comparison with Alternative Strategies}
\label{sec:method:comparison}
% ---------------------------------------------------------------------------

\paragraph{vs.\ \texttt{think@512} only.}
\textsc{Town} achieves substantially higher accuracy at
${\sim}$2.4$\times$ fewer tokens.
The advantage comes from avoiding thinking mode's truncation waste on the
88.8\% of problems where it is unnecessary.

\paragraph{vs.\ \texttt{nothink@256} only.}
\textsc{Town} adds a thinking fallback for the $\sim$11.2\% of hard
problems, recovering accuracy on cases where direct answering
fails.
The additional cost is modest: $\sim$53 extra tokens on average
over the entire dataset (199 vs.\ 146 for \texttt{nothink@256}).

\paragraph{vs.\ Self-consistency~\cite{wang2023selfconsistency}.}
SC@$K$ generates $K$ full-budget reasoning traces and majority-votes,
consuming $K \times B$ tokens unconditionally.
\textsc{Town} uses a \emph{single} non-thinking probe and a binary
routing decision, yielding lower latency (one generation call for easy
questions vs.\ $K$), lower token cost ($\sim$146 tokens for easy
questions vs.\ $K \times B$), and a simpler implementation (no consensus
computation, no threshold tuning).

\paragraph{vs.\ Oracle routing.}
An oracle that knows each problem's difficulty \emph{a priori}
achieves 68.2\% accuracy at 401 average tokens
(Table~\ref{tab:efficiency-frontier}).
\textsc{Town}'s early-stop signal approximates this oracle using a
training-free heuristic, with the gap determined primarily by the
signal's precision on the routed subset.

\paragraph{Hyperparameters.}
\textsc{Town} introduces exactly \emph{two} hyperparameters beyond the
base model: $B_1$ (probe budget) and $B_2$ (thinking fallback budget).
Both have intuitive interpretations: $B_1$ should be large enough for
the model to produce a complete non-thinking answer on most easy
questions (256 tokens suffices for GSM8K), and $B_2$ should be large
enough for a complete reasoning chain on hard questions (512+ for
GSM8K).
No routing thresholds, consensus weights, or learned parameters are
required.
